\subsection{Emotion Geometry: A Story-Derived Space and Content Prompts}
\label{sec:circumplex}

The fifth instrument compares two emotion representations in the Qwen3.5-27B Opus
distill. The first is a story-derived emotion space: 30 emotion centroids per layer,
each a mean over stories the model wrote about a third-person character feeling that
emotion (900 stories in all)~\citep{lyra2026usermodel}. The centroids describe
characters' emotions in generated text, not the state of a conversation partner. The
second is a pair of directions computed from 250 emotional scenario prompts
(hostile$-$calm and desperate$-$calm differences of means), called here the
content-prompt valence and arousal directions. Every number in this subsection is
recomputed from the stored files of~\citet{nexus2026coupling}.

\paragraph{Within the story-derived space.} The space has valence--arousal structure.
From L17 onward its first principal component correlates with the emotions' rated
valence at $|\rho| = 0.66$--$0.80$ and, with the opposite sign, with rated arousal at
$|\rho| = 0.49$--$0.67$, so it runs from pleasant-and-calm to unpleasant-and-activated.
Directions built from the same centroids (hostile$-$calm and desperate$-$calm) align
with that component at $|\cos| = 0.68$--$0.80$ and $0.83$--$0.88$, and the plane
through the hostile, desperate and calm centroids holds $33$--$37\%$ of the centroid
variance. These are within-system figures, so the informative comparison is with
other directions built from the same centroids. From L17, $77$--$89\%$ of the 435
centroid-pair directions align with the first component less than hostile$-$calm does
and $92$--$95\%$ less than desperate$-$calm, and the plane holds $1.26$--$1.36$ times
the variance of the median plane through three centroids.

\paragraph{Across the two prompt sets.} The overlap is weak. At L15, L25, L35 and L45
the content-prompt valence direction has $|\cos| = 0.026$--$0.110$ with the story
space's first component and the arousal direction $0.015$--$0.126$. Scored the same
way, 2,000 random unit directions reach $0.027$--$0.028$ at the 95th percentile and
$0.044$--$0.052$ at most: at L15 the content directions are within that range, and at
L25--L45 they exceed it. Along either content direction lies $0.08$--$0.68\%$ of the
centroid variance, 4 to 34 times a random direction's mean. The model's activations are
not isotropic, so this excess is an upper bound on any shared emotional structure, not
evidence of it.

\paragraph{Across layers.} The two geometries settle on different schedules. The
content prompts' eccentricity is $0.37$--$0.61$ from L17 (it is computed from the norms
of the hostile$-$calm and desperate$-$calm directions, and is 0 for a circle). The
story-derived space's is $0.05$--$0.23$ over L17--L21 and $0.30$--$0.42$ from L22, and
over all 64 layers the two series do not track each other ($r = 0.09$).

\paragraph{Coupling test.} A pre-registered injection experiment added the
content-prompt valence direction, the arousal direction or a random direction at
valence's norm to the output of layer 15, 25, 35 or 45 ($\alpha \in \{1, 3, 5\}$), and
compared the resulting shifts in 250 content prompts and 60 conversational prompts.
The readout is the raw residual stream, which contains the injected vector, so the
cross-pathway cosine is 1.000 at the injection layer by construction. Off that layer
the random direction's mean cosine ($0.577$) is at least as high as valence's ($0.521$)
and arousal's ($0.547$). The experiment shows no direction-specific coupling, and it
cannot test whether the two kinds of processing are separable.
